\documentclass[10pt,a4paper]{amsart}
\usepackage{amsmath,amssymb,amsthm,graphicx}
% Retained minimal preamble for the Stacks Project chapter more-algebra.tex (More on Algebra).
% The chapter file pinned-original-source.tex carries only \input{preamble} on its first line, so no
% preamble text lives inside the chapter itself.  The Stacks Project uses one single shared preamble for
% every chapter; the snapshot stages it identically as sources/stacks-chow/preamble.tex and
% sources/stacks-etale-cohomology/preamble.tex.  The block below is copied mechanically, line for line,
% from sources/stacks-chow/preamble.tex lines 14-18 (the Xy-pic setup) and 154-217 (the native theorem
% declarations together with the mathematical macros).  Nothing was added, removed or redefined; the
% wrapper class and the amsmath/amssymb/amsthm/graphicx base packages are supplied by the assembler.
% Chapter: more-algebra.tex -- More on Algebra.
% For commutative diagrams we use Xy-pic
\usepackage[all]{xy}

% We use 2cell for 2-commutative diagrams.
\xyoption{2cell}

% Theorem environments.
%
\theoremstyle{plain}
\newtheorem{theorem}[subsection]{Theorem}
\newtheorem{proposition}[subsection]{Proposition}
\newtheorem{lemma}[subsection]{Lemma}

\theoremstyle{definition}
\newtheorem{definition}[subsection]{Definition}
\newtheorem{example}[subsection]{Example}
\newtheorem{exercise}[subsection]{Exercise}
\newtheorem{situation}[subsection]{Situation}

\theoremstyle{remark}
\newtheorem{remark}[subsection]{Remark}
\newtheorem{remarks}[subsection]{Remarks}

\numberwithin{equation}{subsection}

% Macros
%
\def\lim{\mathop{\mathrm{lim}}\nolimits}
\def\colim{\mathop{\mathrm{colim}}\nolimits}
\def\Spec{\mathop{\mathrm{Spec}}}
\def\Hom{\mathop{\mathrm{Hom}}\nolimits}
\def\Ext{\mathop{\mathrm{Ext}}\nolimits}
\def\SheafHom{\mathop{\mathcal{H}\!\mathit{om}}\nolimits}
\def\SheafExt{\mathop{\mathcal{E}\!\mathit{xt}}\nolimits}
\def\Sch{\mathit{Sch}}
\def\Mor{\mathop{\mathrm{Mor}}\nolimits}
\def\Ob{\mathop{\mathrm{Ob}}\nolimits}
\def\Sh{\mathop{\mathit{Sh}}\nolimits}
\def\NL{\mathop{N\!L}\nolimits}
\def\CH{\mathop{\mathrm{CH}}\nolimits}
\def\proetale{{pro\text{-}\acute{e}tale}}
\def\etale{{\acute{e}tale}}
\def\QCoh{\mathit{QCoh}}
\def\Ker{\mathop{\mathrm{Ker}}}
\def\Im{\mathop{\mathrm{Im}}}
\def\Coker{\mathop{\mathrm{Coker}}}
\def\Coim{\mathop{\mathrm{Coim}}}

% Boxtimes
%
\DeclareMathSymbol{\boxtimes}{\mathbin}{AMSa}{"02}

%
% Macros for moduli stacks/spaces
%
\def\QCohstack{\mathcal{QC}\!\mathit{oh}}
\def\Cohstack{\mathcal{C}\!\mathit{oh}}
\def\Spacesstack{\mathcal{S}\!\mathit{paces}}
\def\Quotfunctor{\mathrm{Quot}}
\def\Hilbfunctor{\mathrm{Hilb}}
\def\Curvesstack{\mathcal{C}\!\mathit{urves}}
\def\Polarizedstack{\mathcal{P}\!\mathit{olarized}}
\def\Complexesstack{\mathcal{C}\!\mathit{omplexes}}
% \Pic is the operator that assigns to X its picard group, usage \Pic(X)
% \Picardstack_{X/B} denotes the Picard stack of X over B
% \Picardfunctor_{X/B} denotes the Picard functor of X over B
\def\Pic{\mathop{\mathrm{Pic}}\nolimits}
\def\Picardstack{\mathcal{P}\!\mathit{ic}}
\def\Picardfunctor{\mathrm{Pic}}
\def\Deformationcategory{\mathcal{D}\!\mathit{ef}}

% --- appended 2026-10-12 (Chief): the project own definitions of the reference/slogan/history
% display environments, copied mechanically from the shared Stacks preamble lines 8-12
% (\usepackage{verbatim} plus the three \newenvironment lines, which the chapter preamble keeps
% verbatim from the project preamble). Nothing added, removed or redefined: these environments
% carry no mathematical content and the project itself typesets them as comments. Without these
% three lines a counted statement whose own text uses them fails with Environment slogan undefined.
\usepackage{verbatim}
\newenvironment{reference}{\comment}{\endcomment}
\newenvironment{slogan}{\comment}{\endcomment}
\newenvironment{history}{\comment}{\endcomment}
\usepackage{cleveref}
\crefname{originalremark}{original remark}{original remarks}
\Crefname{originalremark}{Original remark}{Original remarks}
\makeatletter
\newlabel{remark-construction}{{09EM}{source}{}{sourcepointer-1}{}}
\newlabel{remark-construction@cref}{{[originalremark][0][]09EM}{[1][0][]source}{}{}{}}
\makeatother
\crefname{originallemma}{original lemma}{original lemmas}
\Crefname{originallemma}{Original lemma}{Original lemmas}
\makeatletter
\newlabel{lemma-extension-dvrs-formally-smooth}{{09E7}{source}{}{sourcepointer-2}{}}
\newlabel{lemma-extension-dvrs-formally-smooth@cref}{{[originallemma][0][]09E7}{[1][0][]source}{}{}{}}
\makeatother
\crefname{originallemma}{original lemma}{original lemmas}
\Crefname{originallemma}{Original lemma}{Original lemmas}
\makeatletter
\newlabel{lemma-pull-root-uniformizer}{{09EV}{source}{}{sourcepointer-3}{}}
\newlabel{lemma-pull-root-uniformizer@cref}{{[originallemma][0][]09EV}{[1][0][]source}{}{}{}}
\makeatother
\crefname{originallemma}{original lemma}{original lemmas}
\Crefname{originallemma}{Original lemma}{Original lemmas}
\makeatletter
\newlabel{lemma-formally-smooth-goes-up}{{09EQ}{source}{}{sourcepointer-4}{}}
\newlabel{lemma-formally-smooth-goes-up@cref}{{[originallemma][0][]09EQ}{[1][0][]source}{}{}{}}
\makeatother
\crefname{originallemma}{original lemma}{original lemmas}
\Crefname{originallemma}{Original lemma}{Original lemmas}
\makeatletter
\newlabel{lemma-multiplicative-e-f}{{0BRL}{source}{}{sourcepointer-5}{}}
\newlabel{lemma-multiplicative-e-f@cref}{{[originallemma][0][]0BRL}{[1][0][]source}{}{}{}}
\makeatother
\title{Abhyankar's lemma: formal smoothness of the extension of the local rings after adjoining roots of the uniformizer}
\author{The Stacks Project Authors}
\date{stacks-project:more-algebra.tex:local-snapshot-20261012; GNU FDL 1.2 or later}
\begin{document}
\maketitle
\setcounter{subsection}{1}
\noindent\textit{Editorial note.} All delivered mathematical text is the author's own, copied from the pinned Stacks Project chapter file more-algebra.tex; no statement or proof step was rewritten, reordered or removed. The counted span is the author environment and its own proof environment. References inside the span are handed to the assembler as navigation stubs that carry the canonical tag of the cited tag only; whether each of them is load-bearing is recorded on the stub itself.\par\medskip
\noindent\textit{Allowed standard prerequisites (editorial).}\begin{itemize}
\item The notation of the construction remark that names the extensions of the discrete valuation rings and the maximal ideals occurring in the statement. used in the statement at line 34041 and in the proof from line 34050. Fixes the objects A\_1, B\_1 and the maximal ideals m\_i, m\_ij so that the statement and the proof talk about the same rings.
\item For an extension of discrete valuation rings, formal smoothness in the adic topology is detected by the ramification index being one and the residue field extension being separable. invoked at proof lines 34070-34072 and again at the end at proof lines 34133-34135. Reduces the two comparison steps to the computation of the ramification index and the separability of the residue field extension.
\item The integral closure of a discrete valuation ring in the extension generated by a root of the uniformizer is a discrete valuation ring with ramification index the order of the root. invoked at proof lines 34092-34095. Identifies the intermediate ring A'\_1 and its ramification index, which lets the argument pass to the subextension generated by the root.
\item Formal smoothness goes up along the diagram of local extensions used to compare the intermediate ring with the original one. invoked at proof lines 34096-34102. Lets the conclusion for the subextension be transferred back to the extension appearing in the statement.
\item Ramification indices of a tower of extensions of discrete valuation rings are multiplicative. invoked at proof lines 34124-34126. Computes the ramification index of the maximal ideal of B\_1 over A\_1 as one in the final case.
\end{itemize}\medskip
\ifdefined\hypertarget\hypertarget{sourcepointer-1}{}\fi
\noindent\textit{Source reference.} Original remark 09EM: sources/stacks-more-algebra/pinned-original-source.tex, source lines 33882--33882. The statement at line 34041 uses the notation of the remark; the remark is a pointer that names the rings and ideals, and every argument in which those objects occur is inside the counted span.\par\medskip
\ifdefined\hypertarget\hypertarget{sourcepointer-2}{}\fi
\noindent\textit{Source reference.} Original lemma 09E7: sources/stacks-more-algebra/pinned-original-source.tex, source lines 32989--32989. Used at proof lines 34070-34072 to reduce the first comparison to a computation of ramification index and residue field, and again at the very end at proof lines 34133-34135; both computations are carried out inside the counted span.\par\medskip
\ifdefined\hypertarget\hypertarget{sourcepointer-3}{}\fi
\noindent\textit{Source reference.} Original lemma 09EV: sources/stacks-more-algebra/pinned-original-source.tex, source lines 33954--33954. Invoked at proof lines 34092-34095 to identify the intermediate ring and its ramification index; the verification that the exponent divides and that the root lies in the larger field is inside the counted span.\par\medskip
\ifdefined\hypertarget\hypertarget{sourcepointer-4}{}\fi
\noindent\textit{Source reference.} Original lemma 09EQ: sources/stacks-more-algebra/pinned-original-source.tex, source lines 34006--34006. Applied at proof lines 34096-34102 to reduce the statement to the case of the subextension generated by the e-th root; constructing that subextension and the diagram it is applied to is inside the counted span.\par\medskip
\ifdefined\hypertarget\hypertarget{sourcepointer-5}{}\fi
\noindent\textit{Source reference.} Original lemma 0BRL: sources/stacks-more-algebra/pinned-original-source.tex, source lines 32952--32952. Used at proof lines 34124-34126 to conclude that the ramification index of the maximal ideal is one; identifying the two extensions making up the tower and the separability of the residue field extension is inside the counted span.\par\medskip
\begin{lemma}[Abhyankar's lemma]
\label{lemma-abhyankar}
Let $A \subset B$ be an extension of discrete valuation rings.
Assume that either the residue characteristic of $A$ is $0$
or it is $p$, the ramification index $e$ is prime to $p$, and
$\kappa_B/\kappa_A$ is a separable field extension.
Let $K_1/K$ be a finite extension. Using the notation of
Remark \ref{remark-construction}
assume $e$ divides the ramification index of $A \subset (A_1)_{\mathfrak m_i}$
for some $i$. Then $(A_1)_{\mathfrak m_i} \subset (B_1)_{\mathfrak m_{ij}}$
is formally smooth in the $\mathfrak m_{ij}$-adic topology
for all $j = 1, \ldots, m_i$.
\end{lemma}

\begin{proof}
Let $\pi \in A$ be a uniformizer. Let $\pi_1$ be a uniformizer
of $(A_1)_{\mathfrak m_i}$. Write $\pi = u \pi_1^{e_1}$ with $u$ a unit
of $(A_1)_{\mathfrak m_i}$ and $e_1$ the ramification index of
$A \subset (A_1)_{\mathfrak m_i}$.

\medskip\noindent
Claim: we may assume that $u$ is an $e$th power in $K_1$.
Namely, let $K_2$ be an extension of $K_1$ obtained by
adjoining a root of $x^e = u$; thus $K_2$ is a factor
of $K_1[x]/(x^e - u)$. Then $K_2/K_1$ is a finite
separable extension (by our assumption on $e$)
and hence $A_1 \subset A_2$ is finite.
Since $(A_1)_{\mathfrak m_i} \to (A_1)_{\mathfrak m_i}[x]/(x^e - u)$
is finite \'etale
(as $e$ is prime to the residue characteristic and $u$ a unit)
we conclude that $(A_2)_{\mathfrak m_i}$ is a factor of
a finite \'etale extension of $(A_1)_{\mathfrak m_i}$ hence
finite \'etale over $(A_1)_{\mathfrak m_i}$ itself.
The same reasoning shows that $B_1 \subset B_2$ induces
finite \'etale extensions
$(B_1)_{\mathfrak m_{ij}} \subset (B_2)_{\mathfrak m_{ij}}$.
Pick a maximal ideal $\mathfrak m'_{ij} \subset B_2$
lying over $\mathfrak m_{ij} \subset B_1$
(of course there may be more than one) and consider
$$
\xymatrix{
(B_1)_{\mathfrak m_{ij}} \ar[r] & (B_2)_{\mathfrak m'_{ij}} \\
(A_1)_{\mathfrak m_i} \ar[u] \ar[r] &
(A_2)_{\mathfrak m'_i} \ar[u]
}
$$
where $\mathfrak m'_i \subset A_2$ is the image.
Now the horizontal arrows have ramification index $1$
and induce finite separable residue field extensions.
Thus, using the equivalence of
Lemma \ref{lemma-extension-dvrs-formally-smooth},
we see that it suffices to show that the right vertical
arrow is formally smooth in the $\mathfrak m'_{ij}$-adic topology.
Since $u$ has a $e$th root
in $K_2$ we obtain the claim.

\medskip\noindent
Assume $u$ has an $e$th root in $K_1$.
Since $e | e_1$ and since $u$ has a $e$th root in $K_1$
we see that $\pi = \theta^e$ for some $\theta \in K_1$.
Let $K'_1 = K[\theta] \subset K_1$ be the subfield generated by $\theta$.
By Lemma \ref{lemma-pull-root-uniformizer} the integral closure $A'_1$
of $A$ in $K[\theta]$ is the discrete valuation ring $A'_1 = A[\theta]$
which has ramification index $e$ over $A$.
If we can prove the lemma for the extension $K'_1/K$,
then we conclude by Lemma \ref{lemma-formally-smooth-goes-up}
applied to the diagram
$$
\xymatrix{
(B'_1)_{B'_1 \cap \mathfrak m_{ij}} \ar[r] & (B_1)_{\mathfrak m_{ij}} \\
A'_1 \ar[u] \ar[r] &
(A_1)_{\mathfrak m_i} \ar[u]
}
$$
for all $j = 1, \ldots, m_i$. This reduces us to the case discussed
in the next paragraph.

\medskip\noindent
Assume $K_1 = K[\pi^{1/e}]$ and set $\theta = \pi^{1/e}$. Let $\pi_B$ be a
uniformizer for $B$ and write $\pi = w \pi_B^e$ for some unit $w$ of $B$.
Then we see that $L_1 = L \otimes_K K_1$ is obtained by adjoining
$\theta / \pi_B$ which is an $e$th root of the unit $w$. Thus
$B \subset B_1$ is finite \'etale. Thus for any maximal ideal
$\mathfrak m \subset B_1$ consider the commutative diagram
$$
\xymatrix{
B \ar[r]_1 & (B_1)_{\mathfrak m} \\
A \ar[u]^e \ar[r]^e & A_1 \ar[u]_{e_\mathfrak m}
}
$$
Here the numbers along the arrows are the ramification indices.
By multiplicativity of ramification indices
(Lemma \ref{lemma-multiplicative-e-f})
we conclude $e_\mathfrak m = 1$. Looking at the residue field extensions
we find that $\kappa(\mathfrak m)$ is a finite separable extension
of $\kappa_B$ which is separable over $\kappa_A$. Therefore
$\kappa(\mathfrak m)$ is separable over $\kappa_A$
which is equal to the residue field of $A_1$ and we win by
Lemma \ref{lemma-extension-dvrs-formally-smooth}.
\end{proof}

\end{document}
